\documentclass[UTF8,fontset=none,12pt]{ctexart}

\setCJKmainfont{Songti SC}
\setCJKsansfont{PingFang SC}
\setCJKmonofont{Songti SC}

\usepackage[a4paper,margin=2.35cm]{geometry}
\usepackage{amsmath,amssymb,amsthm,mathtools,bm}
\usepackage{graphicx}
\usepackage{float}
\usepackage{booktabs}
\usepackage{enumitem}
\usepackage[hidelinks]{hyperref}

\newcommand{\R}{\mathbb{R}}
\newcommand{\epsmach}{\varepsilon_{\mathrm{machine}}}
\newcommand{\fl}{\operatorname{fl}}
\theoremstyle{definition}
\newtheorem{theorem}{定理}
\renewcommand{\thetheorem}{14.\arabic{theorem}}
\renewcommand{\thefigure}{14.\arabic{figure}}

\title{数值线性代数\quad 第 14 讲：稳定性}
\author{Lloyd N. Trefethen、David Bau III\quad 中文阅读版}
\date{}

\begin{document}
\maketitle
\pagestyle{plain}

如果数值算法能够为数值问题提供精确解，那当然是很理想的事情.然而，问题是连续的，而数字计算机是离散的，因此这通常不可能实现.稳定性是刻画“能够做到什么”的标准方式，也是数值分析家理解“得到正确答案”的方式，即使这个答案并不精确.

\section*{算法}

在第 12 讲中，我们把数学问题定义为从数据向量空间 $X$ 到解向量空间 $Y$ 的函数
\[
    f:X\to Y.
\]
算法可以看作同两个空间之间的另一个映射 $\widetilde f:X\to Y$.下面给出这一说法的精确定义.固定一个问题 $f$、一个浮点系统满足式 (13.7)（但不一定满足式 (13.6)）的计算机、一个求解 $f$ 的算法（这里以宽泛意义使用“算法”一词），以及该算法在计算机程序中的实现.

给定数据 $x\in X$，先按照满足式 (13.5) 的方式把数据舍入为浮点数，再将其输入计算机程序.运行程序后，得到一组属于向量空间 $Y$ 的浮点数（因为算法的设计目标就是求解 $f$）.把这个计算结果记为 $\widetilde f(x)$.

这个情形可能非常复杂!至少，$\widetilde f(x)$ 会受到舍入误差的影响.根据具体情况，它还可能受到收敛容差，甚至计算机上同时运行的其他任务等各种因素的影响；例如，当计算分配给处理器的方式直到运行时才确定时，就会出现这种情况.因此，所谓的“函数” $\widetilde f(x)$ 甚至可能在两次运行之间取不同的值；它可能是多值的.（事实上，问题 $f$ 本身也应允许是多值的，这样就能处理允许非唯一解的情形，例如复数的两个平方根都可以接受.）然而，尽管有这些复杂性，我们仍将看到：只依靠式 (13.5) 和 (13.7) 这两个基本公理，就可以对 $\widetilde f(x)$，进而对数值线性代数算法的准确性，作出出人意料地简洁的陈述.

波浪号 $\widetilde{\ }$ 记号非常方便.正如 $\widetilde f$ 是 $f$ 的计算机计算对应物一样，本书中其他计算得到的量也经常用波浪号标记.例如，方程组 $Ax=b$ 的计算解可以记为 $\widetilde x$.

\section*{准确性}

除非是平凡情形，$\widetilde f$ 不可能连续.不过，一个好的算法应当逼近相应的数学问题 $f$.为了定量描述这一点，可以考虑计算的绝对误差
\[
    \|\widetilde f(x)-f(x)\|,
\]
或者相对误差
\begin{equation}
    \frac{\|\widetilde f(x)-f(x)\|}{\|f(x)\|}.
    \tag{14.1}
\end{equation}

本书主要使用相对量，因此式 (14.1) 将作为标准误差度量.

如果 $\widetilde f$ 是一个好的算法，自然会希望相对误差很小，量级为 $\epsmach$.对于每个 $x\in X$，如果算法 $\widetilde f$ 满足
\begin{equation}
    \frac{\|\widetilde f(x)-f(x)\|}{\|f(x)\|}
    =O(\epsmach),
    \tag{14.2}
\end{equation}
就称它对问题 $f$ 是准确的.

粗略地说，式 (14.2) 中的 $O(\epsmach)$ 表示“机器 epsilon 的量级”.不过，$O(\epsmach)$ 也有精确定义，我们马上讨论这一点；这也会说明当分母为零时，类似式 (14.2) 的公式应如何理解.

\section*{稳定性}

如果问题 $f$ 是病态的，那么式 (14.2) 所定义的准确性目标就过于苛刻了.在数字计算机上，输入数据的舍入不可避免；即使后续所有计算都能完美进行，这一扰动本身也可能导致结果发生显著变化.因此，一般情况下最适合追求的不是处处准确，而是稳定性.

如果对每个 $x\in X$，都存在某个 $\widetilde x$，使得下面两式同时成立：
\begin{align}
    \frac{\|\widetilde f(x)-f(\widetilde x)\|}
         {\|f(\widetilde x)\|}
    &=O(\epsmach),
    \tag{14.3}\\[0.5em]
    \frac{\|\widetilde x-x\|}{\|x\|}
    &=O(\epsmach).
    \tag{14.4}
\end{align}
那么称算法 $\widetilde f$ 对问题 $f$ 是稳定的.

用一句话说：
\begin{center}
\textbf{稳定算法对一个几乎正确的问题给出几乎正确的答案.}
\end{center}

这一讲定义的动机将在下一讲以及本书后面的应用中逐渐清楚.需要提醒读者的是：这里给出的稳定性定义在数值线性代数的许多部分都很有用，但条件 $O(\epsmach)$ 对微分方程等其他领域的全部数值问题而言可能过于严格.

\section*{后向稳定性}

许多数值线性代数算法满足一个比稳定性更强、也更简单的条件.如果对每个 $x\in X$，都有
\begin{equation}
    \widetilde f(x)=f(\widetilde x)
    \quad\text{对某个 }\widetilde x\text{ 满足}\quad
    \frac{\|\widetilde x-x\|}{\|x\|}
    =O(\epsmach),
    \tag{14.5}
\end{equation}
则称算法 $\widetilde f$ 对问题 $f$ 是后向稳定的.

这是对稳定性定义的加强：式 (14.3) 中的 $O(\epsmach)$ 被替换成了零.换句话说：
\begin{center}
\textbf{后向稳定算法对一个几乎正确的问题给出完全正确的答案.}
\end{center}

例子将在下一讲给出.

\section*{$O(\epsmach)$ 的含义}

现在解释式 (14.2)--(14.5) 中“$O(\epsmach)$”的精确含义.记号
\begin{equation}
    \varphi(t)=O(\psi(t))
    \tag{14.6}
\end{equation}
是数学中的标准记号，并且有精确定义.它断言存在某个正数 $C$，使得当 $t$ 足够接近某个默认的极限（例如 $t\to0$ 或 $t\to\infty$）时，
\begin{equation}
    |\varphi(t)|\le C\psi(t).
    \tag{14.7}
\end{equation}

例如，陈述 $\sin^2t=O(t^2)$（当 $t\to0$）意味着存在常数 $C$，使得对所有充分小的 $t$，都有 $|\sin^2t|\le Ct^2$.

数学中还有形如
\begin{equation}
    \varphi(s,t)=O(\psi(t))\quad\text{关于 }s\text{ 一致成立}
    \tag{14.8}
\end{equation}
的陈述，其中 $\varphi$ 不仅依赖于 $t$，还依赖于另一个变量 $s$.“一致”表示式 (14.7) 中的同一个常数 $C$ 对 $s$ 的所有取值都成立.例如，
\[
    (\sin^2t)(\sin^2s)=O(t^2)
\]
在 $t\to0$ 时关于 $s$ 一致成立；但如果把正弦换成 $s^2$，这种一致性就会丢失.

本书对符号 $O$ 的使用遵循这些标准定义.具体地，我们经常陈述如下形式的结果：
\begin{equation}
    \|\text{计算得到的量}\|=O(\epsmach).
    \tag{14.9}
\end{equation}

式 (14.9) 的含义如下.首先，“计算得到的量”表示由问题 $f$ 的算法 $\widetilde f$ 所确定的某个数或数的集合的范数，它同时依赖于问题 $f$ 的数据 $x\in X$ 和 $\epsmach$.相对误差 (14.1) 就是一个例子.其次，隐含的极限过程是 $\epsmach\to0$（即 $\epsmach$ 是式 (14.8) 中对应于 $t$ 的变量）.第三，符号 $O$ 对所有数据 $x\in X$ 一致成立（即 $x$ 是对应于 $s$ 的变量）.我们很少专门说明关于 $x\in X$ 的一致性，但它总是隐含的.

在任何一台具体的机器上，$\epsmach$ 都是固定量.谈论极限 $\epsmach\to0$ 时，我们是在考虑计算机，或者更准确地说，计算机族的理想化.式 (14.9) 的意思是：如果把有关算法放到一系列满足式 (13.5) 和 (13.7)、且 $\epsmach$ 逐渐减小到零的计算机上运行，那么 $\|\text{计算得到的量}\|$ 就保证按 $\epsmach$ 的比例或更快地减小.这些理想计算机只要求满足式 (13.5) 和 (13.7)，没有其他要求.

\section*{依赖于 $m,n$，而不是 $A,b$}

我们还可以进一步说明式 (14.2)--(14.5) 中 $O(\epsmach)$ 的含义.隐含常数的一致性可以通过下面的例子说明.假设我们正在研究一个算法，它对非奇异 $m\times m$ 方程组 $Ax=b$ 求解 $x$，并且断言该算法的计算结果 $\widetilde x$ 满足
\begin{equation}
    \frac{\|\widetilde x-x\|}{\|x\|}
    =O\bigl(\kappa(A)\epsmach\bigr).
    \tag{14.10}
\end{equation}

这一断言意味着：存在一个与矩阵 $A$ 和右端 $b$ 都无关的单一常数 $C$，使得对所有充分小的 $\epsmach$，
\begin{equation}
    \frac{\|\widetilde x-x\|}{\|x\|}
    \le C\kappa(A)\epsmach.
    \tag{14.11}
\end{equation}

如果类似式 (14.11) 的公式中分母为零，则采用如下约定.我们真正想表达的是
\begin{equation}
    \|\widetilde x-x\|
    \le C\kappa(A)\epsmach\|x\|.
    \tag{14.12}
\end{equation}

当 $\|x\|\ne0$ 时，两种写法没有区别；但当 $\|x\|=0$ 时，式 (14.12) 清楚地表明，式 (14.10) 的精确含义是：对所有充分小的 $\epsmach$，都有 $\|\widetilde x-x\|=0$.

虽然式 (14.11) 或 (14.12) 中的常数 $C$ 不依赖于 $A$ 或 $b$，但一般会依赖于维数 $m$.从形式上说，这是第 12 讲问题定义的结果.如果定义问题 $f$ 的维数（例如 $m$ 或 $n$）发生变化，那么向量空间 $X$ 和 $Y$ 也必须变化，于是我们实际上得到了一个新问题 $f'$.

从实际角度看，数值线性代数算法中的舍入误差影响通常也会随 $m$ 和 $n$ 增长.不过，这种增长通常足够慢，不会造成严重问题.在最坏情况下，对 $m$ 或 $n$ 的依赖通常是线性、二次或三次的（指数取决于所选范数以及所选算法）；而由于统计抵消，大多数数据的误差远小于最坏情况.

原则上，类似式 (14.9) 的陈述可能隐藏一个 $2^m$ 这样的维数相关因子，使界在实践中毫无用处.然而，本书中只有一个地方会出现这种情况，即带部分选主元的高斯消去法；届时我们会充分提醒读者，避免误解.一般而言，当本书印出 $O(\epsmach)$ 时，在真实机器上进行真实计算时，相关量通常至多是 $\epsmach$ 的 100 倍，或者也许 1000 倍.

\section*{与范数的独立性}

涉及 $O(\epsmach)$ 的定义有一个方便的性质：只要 $X$ 和 $Y$ 是有限维空间，这些定义就与范数的选择无关.

\begin{theorem}
对于定义在有限维空间 $X$ 和 $Y$ 上的问题 $f$ 及算法 $\widetilde f$，准确性、稳定性和后向稳定性这三个性质是否成立，都独立于 $X$ 和 $Y$ 中范数的选择.
\end{theorem}

\textbf{证明.}众所周知（也很容易证明），在有限维向量空间中，所有范数都是等价的.也就是说，如果 $\|\cdot\|$ 和 $\|\cdot\|'$ 是同一空间上的两个范数，那么存在正数 $C_1,C_2$，使得对该空间中的所有 $x$，
\[
    C_1\|x\|\le\|x\|'\le C_2\|x\|.
\]
因此，改变范数可能改变 $O(\epsmach)$ 陈述中隐含常数 $C$ 的大小，但不会改变这样一个常数是否存在. \hfill$\square$

\section*{习题}

\begin{enumerate}[label=\textbf{14.\arabic*}]
    \item 判断正误.
    \begin{enumerate}[label=(\alph*)]
        \item $\sin x=O(1)$，当 $x\to\infty$.
        \item $\sin x=O(1)$，当 $x\to0$.
        \item $\log x=O(x^{1/100})$，当 $x\to\infty$.
        \item $n!=O((n/e)^n)$，当 $n\to\infty$.
        \item $A=O(V^{2/3})$，当 $V\to\infty$；其中 $A$ 和 $V$ 分别是球的表面积和体积，度量单位分别为平方英里和立方微米.
        \item $\fl(\pi)-\pi=O(\epsmach)$.（这里没有写出极限 $\epsmach\to0$，因为本书所有 $O(\epsmach)$ 表达式都默认这一点.）
        \item $\fl(n\pi)-n\pi=O(\epsmach)$，对所有整数 $n$ 一致成立.（这里的 $n\pi$ 表示精确的数学量，而不是浮点计算的结果.）
    \end{enumerate}

    \item
    \begin{enumerate}[label=(\alph*)]
        \item 证明
        \[
            (1+O(\epsmach))(1+O(\epsmach))
            =1+O(\epsmach).
        \]
        这一陈述的精确含义是：如果函数 $f$ 满足
        \[
            f(\epsmach)
            =(1+O(\epsmach))(1+O(\epsmach))
            \quad\text{当 }\epsmach\to0,
        \]
        那么也有
        \[
            f(\epsmach)=1+O(\epsmach)
            \quad\text{当 }\epsmach\to0.
        \]
        \item 证明
        \[
            (1+O(\epsmach))^{-1}=1+O(\epsmach).
        \]
    \end{enumerate}
\end{enumerate}

\end{document}
